\documentclass[11pt,reqno,a4paper]{amsart}

\usepackage[T1]{fontenc}
\usepackage{lmodern}
\usepackage{amsmath,amssymb,mathtools}
\usepackage{array}
\usepackage{booktabs}
\usepackage{enumitem}
\usepackage[protrusion=true,expansion=false]{microtype}
\usepackage{xcolor}
\usepackage[hidelinks]{hyperref}

\newtheorem{theorem}{Theorem}[section]
\newtheorem{proposition}[theorem]{Proposition}
\newtheorem{lemma}[theorem]{Lemma}
\newtheorem{corollary}[theorem]{Corollary}
\theoremstyle{definition}
\newtheorem{definition}[theorem]{Definition}
\theoremstyle{remark}
\newtheorem{remark}[theorem]{Remark}

\newcommand{\Fin}{\operatorname{Fin}}
\newcommand{\ProblemAt}{\operatorname{Problem617At}}
\newcommand{\IndepFree}{\operatorname{IndepSetFree}}
\newcommand{\CliqueFree}{\operatorname{CliqueFree}}
\newcommand{\Admissible}{\operatorname{Admissible}}
\makeatletter
\newcommand{\sha}[1]{%
  \begingroup\ttfamily\small
  \@tfor\@tempa:=#1\do{\@tempa\hspace{0pt}}%
  \endgroup}
\makeatother
\newcommand{\code}[1]{\texttt{#1}}

\title[Machine-verified fixed \(r=5\) case]
{Machine verification of the fixed \(r=5\) case of Erd\H{o}s Problem 617}
\author{Ramazan Kara}
\date{24 July 2026}
\subjclass[2020]{05C15, 05C35, 68V20, 03B35}
\keywords{edge coloring, formal verification, Lean, LRAT, Erd\H{o}s problem}

\begin{document}
\hypersetup{
  pdftitle={Machine verification of the fixed r=5 case of Erdos Problem 617},
  pdfauthor={Ramazan Kara}
}

\begin{abstract}
We give an independent machine verification of the fixed \(r=5\) case of an
Erd\H{o}s--Gy\'arf\'as balanced-coloring conjecture.  In exact form, the Lean
theorem states that every five-coloring of the unordered edges of \(K_{26}\)
has a six-vertex set on which some color is absent.  The mathematical source
is Robert Sneiderman's 2026 preprint.  The contribution reported here is an
independent formal and certificate verification: the coloring semantics,
graph reductions, extremal endpoint arguments, residual contradiction, and
the final composition are checked in Lean 4.32.0 with a pinned mathlib
revision.  The remaining finite special-Brooks obstruction is split into 89
exhaustive branches.  Each branch has a RUP-only LRAT certificate checked by
two external checkers and imported as a propositional theorem into Lean; a
kernel-checked semantic and symmetry bridge connects these theorems to the
unrestricted graph proposition.  A fresh audit checks 89 kernel imports, 89
graph-semantic closures, complete branch coverage, and 192 assumption
queries.  The final theorem uses only \(\mathsf{propext}\),
\(\mathsf{Classical.choice}\), and \(\mathsf{Quot.sound}\).

This verifies only the fixed \(r=5\) case.  It does not settle the
conjecture for all \(r\geq3\), which remains open.
\end{abstract}

\maketitle

\begin{center}
\fbox{\parbox{0.91\textwidth}{\centering
\textbf{Scope notice.}
The machine-verified theorem in this paper is \(\ProblemAt(5)\) only.
The unrestricted all-\(r\) Erd\H{o}s Problem 617 conjecture remains open.
No claim of independent external review is made.}}
\end{center}

\medskip
\noindent\textbf{Attribution and provenance.}
The original mathematical proof is due to Robert Sneiderman and appears in
the Robby955 preprint repository at pinned commit
\sha{735c2eeffcb5c8abe92f8a6be04be2fa3c8bb6ab}
\cite{Sneiderman2026}.  This paper describes a separate formal-verification
project.  The formalization is not represented as authorship of the original
mathematical proof, nor as external review of that preprint.

\section{The fixed theorem and its boundary}

Erd\H{o}s and Gy\'arf\'as asked whether the following assertion holds for
every integer \(r\geq3\): whenever the edges of \(K_{r^2+1}\) are colored
with \(r\) colors, some \(r+1\) vertices induce a complete graph on which at
least one color is absent \cite{ErdosGyarfas1999,Erdos1999}.  The question is
listed as Erd\H{o}s Problem 617 \cite{ErdosProblems617}.

For finite types \(V\) and \(C\), an edge coloring is a total function on the
unordered, non-diagonal pairs of vertices,
\[
  \chi:\binom V2\longrightarrow C.
\]
For \(S\subseteq V\), a color \(c\in C\) is absent on \(S\) when
\(\chi(e)\ne c\) for every \(e\in\binom S2\).  The fixed statement checked in
this project is the following.

\begin{theorem}[Fixed \(r=5\) upper theorem]\label{thm:r5}
For every map
\[
  \chi:\binom{[26]}2\longrightarrow[5],
\]
there are \(S\in\binom{[26]}6\) and \(c\in[5]\) such that \(c\) is absent
from every edge induced by \(S\).
\end{theorem}

The corresponding lower bound on 25 vertices is the standard affine
construction.  Take \(V=\mathbb F_5^2\).  The six parallel classes of affine
lines are indexed by the five slopes and the vertical direction.  Merge two
classes and use the other four as individual colors.  Every six points
contain two points on a common line in each of the five resulting color
classes.  Thus every six-set sees every color on \(K_{25}\).  Together with
Theorem~\ref{thm:r5}, this gives the equivalent set-coloring Ramsey identity
\(R(6;5,4)=26\), as explained in \cite{Sneiderman2026}.

\begin{remark}[Why this is not a solution of the full problem]
The formal development defines
\[
  \ProblemAt(r) :=
  \text{the upper assertion on }K_{r^2+1}\text{ with }r\text{ colors},
\]
and separately defines
\[
  \operatorname{Problem617}:=
  \forall r,\ 3\le r\longrightarrow\ProblemAt(r).
\]
The exported theorem has type \(\ProblemAt(5)\), not
\(\operatorname{Problem617}\).  A positive result for one fixed value cannot
be generalized by changing a numeral.
\end{remark}

\section{Exact Lean statement}

The formal model uses
\code{SimpleGraph.TopEdgeLabeling}.  This prevents two common specification
errors: loops are not silently colored, and the two orientations of an edge
cannot receive different colors.  The central definitions, with notation
slightly compressed, are:
\begin{align*}
  \operatorname{ColorPresent}(\chi,S,c)
    &:\!\!\iff \exists u,v\in S,\ u\ne v\ \wedge\ \chi(\{u,v\})=c,\\
  \operatorname{ColorAbsent}(\chi,S,c)
    &:\!\!\iff \neg\operatorname{ColorPresent}(\chi,S,c),\\
  \operatorname{UpperStatement}(V,C,k)
    &:\!\!\iff \forall\chi,\ \exists S,\ |S|=k\ \wedge{}\\[-2pt]
    &\hspace{5.2em}\exists c,\ \operatorname{ColorAbsent}(\chi,S,c).
\end{align*}
The fixed theorem is definitionally
\[
 \code{R5Upper}
 =\operatorname{UpperStatement}(\Fin(26),\Fin(5),6),
\]
while
\[
 \code{Problem617At}\ r
 =\operatorname{UpperStatement}(\Fin(r^2+1),\Fin(r),r+1).
\]
Lean checks by reflexivity that
\(\code{Problem617At}\ 5\leftrightarrow\code{R5Upper}\).

For each color \(c\), the color graph \(G_c\) contains exactly the edges
colored \(c\).  The development proves
\[
 \operatorname{ColorAbsent}(\chi,S,c)
 \ \longleftrightarrow\
 G_c.\operatorname{IsIndepSet}(S),
\]
and consequently
\[
 \operatorname{IsCounterexample}(6,\chi)
 \ \longleftrightarrow\
 \forall c,\ G_c.\IndepFree(6).
\]
These are ordinary propositions proved from the definitions; the
certificate layer does not redefine the original coloring statement.

The final declarations are:
\begin{center}
\begin{tabular}{@{}ll@{}}
\toprule
Lean declaration & Type\\
\midrule
\code{e058R5SpecialBrooksObstruction}
  & \code{R5SpecialBrooksObstruction}\\
\code{e058NoR5Counterexample}
  & \(\neg\exists\chi,\ \code{IsCounterexample 6 }\chi\)\\
\code{e058R5Upper}
  & \code{R5Upper}\\
\code{e058Problem617AtFive}
  & \code{Problem617At 5}\\
\bottomrule
\end{tabular}
\end{center}

\section{Mathematical architecture}

\subsection{Color graphs and admissibility}

Assume for contradiction that every six vertices see all five colors.
For a fixed color graph \(G\), every six-set \(S\) satisfies
\begin{equation}\label{eq:local}
  1\le e(G[S])\le11.
\end{equation}
The lower bound says that the fixed color occurs.  The upper bound follows
because each of the other four colors occupies at least one of the 15
edges.  In particular,
\[
  \alpha(G)\le5,\qquad \omega(G)\le5.
\]
The formal development calls a graph \emph{admissible} when every six
vertices span at most 11 edges.  Admissibility is inherited by induced
subgraphs.

The five color graphs partition \(E(K_{26})\), so a least color has at most
\(325/5=65\) edges.  The mathematical proof uses Brooks's theorem and
specialized extremal estimates to confine its minimum degree
\cite{Brooks1941,KangPikhurko2005,Sneiderman2026}.  In the formal
development, the imported general theorems are replaced by direct checked
specializations.  The only finite endpoint isolated as a proposition is
\begin{equation}\label{eq:brooks}
\begin{gathered}
\operatorname{R5SpecialBrooksObstruction}:\\
\nexists\,G\text{ on }\Fin(26)\text{ that is 5-regular and admissible,}\\
G.\CliqueFree(6)\quad\text{and}\quad G.\IndepFree(6).
\end{gathered}
\end{equation}
Sections~\ref{sec:certificates} and \ref{sec:bridge} explain how
\eqref{eq:brooks} is proved in the kernel rather than assumed.

\subsection{Independence layers}

The non-certificate portion of the formal proof develops three layers of
induced-subgraph estimates.  We record their roles because they explain the
later residual recursion.

\begin{enumerate}[leftmargin=2.2em]
\item If \(F\) is admissible, \(\alpha(F)\le2\), and \(n=|V(F)|\ge12\),
then the complement has maximum degree at most four, giving
\[
  e(F)\ge\binom n2-2n.
\]
The order-11 endpoint is proved separately: \(e(F)\ge36\).

\item For admissible graphs with \(\alpha(F)\le3\), the formal proof
propagates the first layer and proves exact lower bounds 45, 50, and 54 at
orders 16, 17, and 18.  The order-10 equality analysis identifies the
balanced two-fold blow-up of \(C_5\) and its exact cover structure.

\item For admissible graphs with \(\alpha(F)\le4\), the next propagation
gives lower bounds 55, 59, and 62 at orders 21, 22, and 23.  Exact
order-15 structures at 35 and 36 edges supply the equality cases required
by the final argument.
\end{enumerate}

These statements are quantified over arbitrary finite simple graphs with the
displayed hypotheses.  The small transparent enumerations use kernel
\code{decide}, not \code{native\_decide}; no graph catalog or solver result is
an unmentioned premise.

\subsection{Equalization and residual contradiction}

Conditional on \eqref{eq:brooks}, the least-color reduction yields a vertex
of degree two, three, or four.  The independence-layer bounds exclude
degrees two and three.  They also force every color graph to have exactly 65
edges.  Equality in the neighborhood accounting isolates a \(K_5\), leaving
an induced graph \(H\) on 21 vertices with
\[
  e(H)=55,\qquad \alpha(H)\le4,\qquad
  \delta(H)\ge4,
\]
and with inherited admissibility.

The degree-four residual branch peels a second isolated \(K_5\) and uses the
order-10 and order-11 endpoints to prove \(\delta(H)\ge5\).  The degree-five
branch uses the exact order-15 alternatives, common-avoider lemmas, cover
transport through the canonical two-fold \(C_5\) blow-up, and the local
six-set cap to derive a contradiction.  Lean therefore proves
\[
  \operatorname{R5SpecialBrooksObstruction}
  \longrightarrow
  \neg\exists\chi,\operatorname{IsCounterexample}(6,\chi),
\]
and composes this implication with the exact coloring semantics.

The proof-module boundary is summarized in
Table~\ref{tab:formal-layers}.
\begin{table}[ht]
\centering
\caption{Principal non-certificate formal layers.}
\label{tab:formal-layers}
\footnotesize
\begin{tabular}{@{}
  >{\raggedright\arraybackslash}p{0.18\textwidth}
  >{\raggedright\arraybackslash}p{0.37\textwidth}
  >{\raggedright\arraybackslash}p{0.37\textwidth}@{}}
\toprule
Stage & Lean module & Checked role\\
\midrule
Semantics & \code{Basic} & Exact coloring and color-graph equivalence\\
Large terminals & \code{AlphaTwoLarge}, \code{AlphaPropagation}
  & Direct admissible density bounds\\
Exact endpoints & \code{AlphaTwoEleven}, \code{OrderTenStructure}
  & Orders 10 and 11\\
Middle layers & \code{AlphaThreeSmall}, \code{OrderFifteenStructure}
  & Independence three and exact order 15\\
Last layer & \code{AlphaFourSmall}
  & Independence four\\
Equalization & \code{EdgeEqualization}
  & 65 edges and first isolated \(K_5\)\\
Residuals & \code{ResidualDegreeFour}, \code{ResidualDegreeFive}
  & Final contradiction, conditional on \eqref{eq:brooks}\\
\bottomrule
\end{tabular}
\end{table}

\section{The 89 certificate branches}\label{sec:certificates}

\subsection{Exhaustive finite split}

Fix a graph satisfying the four hypotheses of
\eqref{eq:brooks}.  Relabel a selected vertex as 0 and its five neighbors as
\(1,\ldots,5\).  There are 26 canonical isomorphism types for the graph
induced by those five neighbors.  The remaining 20 vertices are sorted by
their five-bit adjacency rows.  The difficult neighborhood type, branch 19,
is refined by cross patterns and then by a zero-pattern anchor.  The
complete leaf inventory is:
\begin{center}
\begin{tabular}{@{}lrl@{}}
\toprule
Family & Leaves & Meaning\\
\midrule
E038 & 19 & Neighborhood branches \(0,\ldots,18\)\\
E042 & 17 & Direct cross-pattern children of branch 19\\
E043 & 47 & Residual zero-anchor children: \(26+16+5\)\\
E045 & 6 & Final quotient branches \(20,\ldots,25\)\\
\midrule
Total & 89 & Pairwise named exhaustive leaves\\
\bottomrule
\end{tabular}
\end{center}

Each leaf is encoded as a CNF over 325 primary graph-edge variables and
deterministic auxiliary variables for cardinality and pattern constraints.
The E038, E042, and E043 leaves share an allocation scheme with 3930
variables; the six E045 leaves use a second scheme with 3920.  The
clauses follow necessarily from regularity of degree five, admissibility,
the absence of a six-clique or an independent six-set, and the branch's
canonical unit assignment.  Hence unsatisfiability of every leaf excludes the
graph family, provided both encoding soundness and branch coverage are proved.

\subsection{RUP-core slicing}

The branch proofs are LRAT traces containing only reverse unit-propagation
additions.  A deterministic backward slicer starts at the final
empty clause, follows every antecedent, retains precisely the reachable
original and derived clauses, and renumbers them monotonically.  It rejects
RAT hints, malformed or repeated identifiers, missing antecedents, comments
in the kernel input, and a missing final empty clause.

For every leaf, the reduced CNF/LRAT pair is accepted by:
\begin{enumerate}[leftmargin=2.2em]
\item the compiled C checker derived from \code{drat-trim}
      \cite{WetzlerHeuleHunt2014}; and
\item a separately implemented Python RUP checker.
\end{enumerate}
Both checkers reject two deterministic corruptions per leaf: a truncated
proof and a forged empty clause.  This gives four negative checker outcomes
per leaf, so the replay ledger contains 89 C checks, 89 Python checks, and 356
negative checks.  Separate unit tests reject dependency and renumbering
damage.  These external checks are defense in depth.  They do not replace the
Lean import.

\begin{table}[ht]
\centering
\caption{Aggregate original and backward-core certificate inventory.}
\label{tab:cores}
\begin{tabular}{@{}lrr@{}}
\toprule
Artifact family & Bytes & Clauses/additions\\
\midrule
Original CNFs (decompressed)
  & \(2{,}709{,}936{,}427\) & \(43{,}878{,}934\) clauses\\
Original LRATs (decompressed)
  & \(795{,}950{,}803\) & \(1{,}274{,}831\) additions\\
Reduced CNFs
  & \(5{,}723{,}342\) & \(183{,}145\) clauses\\
Reduced LRATs
  & \(441{,}983{,}911\) & \(1{,}274{,}831\) additions\\
Old-to-new maps
  & \(48{,}467{,}568\) & 89 maps\\
\bottomrule
\end{tabular}
\end{table}
Every one of the \(1{,}274{,}831\) retained additions is RUP; the aggregate
RAT count is zero.  The slicing gain comes chiefly from removing unused
original clauses and deletion records while retaining the complete reachable
derivation.

\subsection{Lean LRAT import}

Mathlib's LRAT machinery constructs a term of type
\[
  \code{Sat.Fmla.proof}\ \Gamma\ [\,],
\]
which says that every valuation satisfying the retained input formula
\(\Gamma\) satisfies the empty clause \cite{CruzFilipeEtAl2017,mathlib2020}.
Large proofs are split into dependency-ordered stages to bound elaborator
memory.  Each stage imports the theorem produced by the previous stage; the
last stage exports a named unsatisfiability theorem.  The Lean kernel checks
all stages.  The generated theorem sources and final object files are hashed
in per-leaf receipts.

Backward slicing can delete primary input clauses that are unnecessary for a
particular refutation.  For this reason, a hashed unit manifest is generated
from the exact original branch formulas, not inferred from the reduced core.
The semantic layer proves that any graph satisfying the branch hypotheses
gives a valuation satisfying:
\begin{enumerate}[leftmargin=2.2em]
\item every retained primary graph clause;
\item every retained sequential-counter clause;
\item every retained exterior-pattern clause; and
\item the exact canonical unit assignment.
\end{enumerate}
Combining this valuation theorem with the imported empty-clause theorem gives
one graph-semantic contradiction per leaf.

\section{Kernel-checked symmetry and semantic bridge}\label{sec:bridge}

Certificate checking alone would prove only that 89 propositional formulas
are unsatisfiable.  E058 therefore includes a separate proof of the route
from an arbitrary graph in \eqref{eq:brooks} to one of those formulas.

\subsection{Normalization}

The development constructs an explicit vertex permutation that fixes vertex
0 and sends its actual five-element neighborhood to
\(\{1,2,3,4,5\}\).  It proves that graph relabeling preserves degree,
admissibility, clique-freeness, and independence-freeness.  A second
permutation fixes the first six vertices and lexicographically sorts the 20
exterior rows.

\subsection{Orbit coverage}

Transparent generated witness tables prove that every five-vertex
neighborhood graph belongs to one of the 26 canonical orbits.  Branch 19
receives a separately checked cross-pattern orbit table.  The three remaining
zero-anchor parents receive checked column and row actions and witness-child
tables.  These computations use ordinary reducible definitions and kernel
\code{decide}; their generator output is regenerated byte-for-byte in the
fresh audit.

The coverage modules establish:
\begin{align*}
 &\text{E038 leaves}\ \Longrightarrow\ \text{branches }0,\ldots,18
   \text{ impossible},\\
 &\text{E042/E043 leaves}\ \Longrightarrow\ \text{branch }19
   \text{ impossible},\\
 &\text{E045 leaves}\ \Longrightarrow\ \text{branches }20,\ldots,25
   \text{ impossible}.
\end{align*}
The final composition is the following checked implication chain:
\[
\begin{array}{c}
\text{89 kernel LRAT theorems}\\
\Downarrow\\
\text{89 graph-semantic contradictions}\\
\Downarrow\\
\text{26 canonical neighborhood branches impossible}\\
\Downarrow\\
\code{R5SpecialBrooksObstruction}\\
\Downarrow\\
\code{R5Upper}\ \Longleftrightarrow\ \code{Problem617At 5}.
\end{array}
\]
No symmetry assumption is inserted as an axiom, and no receipt or Python
Boolean is coerced into a theorem.

\section{Fresh audit and verification level}\label{sec:audit}

The recorded final release validation is one uninterrupted audit of exact
commit
\sha{d19a0cf786a0fa714289830f276cf406408ab65b}.  The mirror contained the
separately compiled C checker but no local Lean build directory when the
audit began.  The process completed every gate after 25,497.97 seconds.  Its
receipt records zero warnings and zero forbidden source hits.  Its
\code{fresh\_local\_build} field is true, and its
\code{resumed\_after\_core\_replay} field is false.  The receipt's SHA-256 is
\sha{789acb2bc829fcba0a5656fabca5a37d9d4fdd2b493858039ad0081cc37032d1}.

Earlier development evidence remains retained with its original
qualifications.  A successful clean-start sequence used post-core resume
segments, and a later hardened final-source phase explicitly reported itself
as a resume while revalidating retained stages, generated sources, coverage
modules, and assumptions.  A provenance receipt binds that history to its
legacy runner and terminal receipt.  Two later exact-source attempts that
produced no receipt are also retained as infrastructure diagnostics: one was
terminated at a command-session boundary, and one lacked the separately built
C checker.  The uninterrupted exact-commit run removes the resume
qualification from the final-source evidence without rewriting this history.

All successful runs pin Lean to \code{leanprover/lean4:v4.32.0} and mathlib to revision
\sha{81a5d257c8e410db227a6665ed08f64fea08e997}.  The audit then performs the
following gates in order:
\begin{enumerate}[leftmargin=2.2em]
\item regenerate the three transparent orbit-certificate sources and compare
      them byte-for-byte;
\item build the base library and the E058 semantic bridge with zero warnings;
\item regenerate all 89 reduced RUP cores from the committed compressed
      inputs;
\item run both external checkers and all 356 corruption tests;
\item regenerate and hash all 89 primary unit CNFs;
\item re-split every core, regenerate each committed stage source
      byte-for-byte, and freshly import every staged kernel proof;
\item build all 89 graph-semantic closure theorems;
\item compile all eight coverage modules;
\item query the assumptions of 89 kernel theorems, 89 semantic theorems, and
      14 coverage/final theorems; and
\item scan every audited source for \code{sorry}, \code{admit}, project
      \code{axiom}, \code{unsafe}, \code{native\_decide},
      \code{ofReduceBool}, and \code{run\_tac}; compare source hashes before
      and after the run.
\end{enumerate}

\begin{table}[ht]
\centering
\caption{Fresh E058 audit ledger.}
\label{tab:audit}
\begin{tabular}{@{}lr@{}}
\toprule
Gate & Passed items\\
\midrule
Reduced RUP cores & 89\\
External C/Python positive checks & \(89+89\)\\
Deliberate corruption rejections & 356\\
Lean kernel LRAT imports & 89\\
Lean graph-semantic closures & 89\\
Coverage modules & 8\\
Assumption queries & 192\\
Warnings & 0\\
Forbidden source hits & 0\\
\bottomrule
\end{tabular}
\end{table}

\paragraph{Diagnostic benchmark outcome.}
The preregistered proof-only import of the unreduced E038 branch-02
CNF/LRAT did not complete within its 300-second cap on the final audit
machine.  GNU \code{time} recorded exit status 124, 310.08 user seconds,
7.57 system seconds, 5:19.34 elapsed including termination grace and cleanup,
and peak resident memory of 5,924,372~KiB.  The exact source, empty Lean
output, and resource log are retained.  Consequently no raw-versus-core
speedup ratio is claimed.  The theorem-bearing reduced core was independently
regenerated, checked, staged, and accepted by Lean in every successful E058
audit.  This failed diagnostic performance gate is not a logical premise of
the kernel theorem, but it is reported because it was preregistered.

The assumption audit reports that
\code{Erdos617.e058Problem617AtFive} depends exactly on
\[
 \{\mathsf{propext},\mathsf{Classical.choice},\mathsf{Quot.sound}\}.
\]
These are standard logical foundations used by mathlib.  In particular, the
dependency list contains no \code{sorryAx} and no project-defined axiom.

Under the repository's verification vocabulary, this is a
\emph{machine-verified resolution of the fixed \(r=5\) scope}.  It is not an
externally verified result: the formal development and this preprint have not
yet completed independent expert review.  It is also not a resolution of the
quantified all-\(r\) conjecture.

\section{Disclosure of AI assistance}

OpenAI Codex was used as an AI-assisted research and software engineering
tool in this project.  It helped draft and revise Lean, Python, shell,
documentation, and \LaTeX{} material; execute verification workflows; inspect
intermediate results; and organize the reproducibility record.  The formal
and expository output should therefore not be read as unaided human work.
Codex is not an author and its textual assertions are not evidence for any
mathematical claim.  The internal adversarial audit was likewise
Codex-assisted and is not represented as independent human review.

Every machine-verification claim in this paper is instead tied to committed
review source, certificates, kernel-checked theorem objects, audit
transcripts, and cryptographic hashes that can be replayed independently.
Ramazan Kara directed the project, selects what may be published, and assumes
responsibility for any published artifact and for this account.  Independent
expert review has not yet been completed.

\section{Trust boundary}

It is useful to separate proof premises from reproducibility checks.

\paragraph{Logically trusted.}
The final theorem is accepted by the Lean 4 kernel.  The practical trusted
computing base includes the Lean parser, elaborator and kernel executable,
the operating system and hardware used to run it, and the standard logical
axioms printed above \cite{deMouraUllrich2021}.  The pinned mathlib sources
supply graph theory infrastructure and the LRAT proof constructor
\cite{mathlib2020}.

\paragraph{Not trusted as theorem oracles.}
The SAT solvers that originally produced the LRAT traces are not theorem
premises.  Neither are the Python and C external checkers, the core slicer,
the certificate generators, the audit receipts, or hash comparisons.  An
error in one of these tools can at worst cause it to propose a bad proof term
or a mismatched source; the Lean kernel must still accept the resulting term,
and the graph-semantic and coverage theorems must still type-check.

\paragraph{Why hashes still matter.}
Kernel checking establishes the proposition compiled in a particular run.
Hashes establish that the reviewed or eventually released sources, compressed
certificates, generated units, transcripts, and receipts are the same
artifacts described by this paper.  They close the distribution and replay
boundary, not the logical inference boundary.

\paragraph{Excluded claims.}
No assertion is made here about the correctness of claimed fixed cases
\(r=6,\ldots,9\), except that they are prior work outside this verification.
No unsuccessful search, solver agreement, or unrestricted SAT timeout is
used as evidence.  The full statement
\(\forall r\ge3,\ProblemAt(r)\) remains open.

\section{Reproduction}\label{sec:reproduction}

The local release candidate contains:
\begin{itemize}[leftmargin=2em]
\item all pinned Lean source files, \code{lean-toolchain},
      \code{lake-manifest.json}, and the project lakefile;
\item the 89 original compressed CNFs and LRAT certificates;
\item the core slicer, both external checkers, stage generator, semantic
      runner, complete fresh-audit runner, and their unit tests;
\item the exact unit manifest, fresh receipts, audit source and output,
      transcripts, resource records, and SHA-256 inventory;
\item this complete \LaTeX{} source, bibliography, build file, and rendered
      PDF.
\end{itemize}

From the repository root, build the small external checker and run the
complete E058 unit-test target:
\begin{verbatim}
make verify-e058-unit
\end{verbatim}
Create a source mirror that has no \code{formal/lean/.lake/build} directory
and run:
\begin{verbatim}
python3 repro/run_e058_special_brooks_kernel_audit.py \
  --project /fresh --lake /lean-4.32.0/bin/lake \
  --log /fresh/e058-audit.log --require-no-local-build
\end{verbatim}
The success line begins
\begin{verbatim}
E058-FRESH-LEAN-AUDIT-PASS cores=89 kernel_imports=89
semantic_closures=89 axiom_queries=192
warnings=0 forbidden_hits=0
\end{verbatim}
The retained final audit ran this command without interruption from exact
commit \sha{d19a0cf786a0fa714289830f276cf406408ab65b}; the C checker was
compiled first, while the local Lean build directory remained absent.  Its
receipt has SHA-256
\sha{789acb2bc829fcba0a5656fabca5a37d9d4fdd2b493858039ad0081cc37032d1}.
The earlier clean-start/resume sequence, hardened replay, receipt-assembly
failure, and two exact-source infrastructure failures remain retained and
explicitly labeled; none is substituted for the uninterrupted receipt.
Finally, verify the release inventory with:
\begin{verbatim}
sha256sum --check \
  artifacts/e058_special_brooks_kernel_bridge/SHA256SUMS
\end{verbatim}
The release reproduction guide records the exact mirror command, expected
tool hashes, receipt hashes, and required disk and memory envelope.

\section{Conclusion}

The checked declaration
\[
  \code{Erdos617.e058Problem617AtFive} :
  \code{Problem617At 5}
\]
establishes that every five-coloring of \(K_{26}\) has six vertices omitting
a color.  The result is an independent formal and certificate verification
of the fixed-\(r=5\) theorem whose mathematical source is Sneiderman's
preprint.  Its main verification contribution is the end-to-end bridge from
89 LRAT certificates, through exact graph semantics and exhaustive symmetry
coverage, to the unconditional original coloring statement.

The quantified Erd\H{o}s Problem 617 conjecture for every \(r\ge3\) remains
open.

\bibliographystyle{alpha}
\bibliography{references}

\end{document}
